\documentclass[../../main.tex]{subfiles}

\begin{document}

\section{Fundamental Properties}

Let's first start by solving some problems.

\begin{problem}{}{}
Identify if the set $\mathbb{S} = \{ (x, y) \in \mathbb{R}^2 \mid
x + y = 5 \}$ is a vector space.
\end{problem}
\begin{solution}
We could see if the set satisfies all axioms for a vector space.
However, just by looking at the condition $x + y = 5$, we know that
there cannot exist a zero vector since $0 + 0 \neq 5$. Because the
set violates A4, the set is not a vector space.
\end{solution}

\begin{problem}{}{}
Identify if the set $\mathbb{S} = \{ \langle x, y \rangle \in
\mathbb{R}^2 \mid x + y = 0 \}$ is a vector space.
\end{problem}
\begin{solution}
Unlike the previous problem, we know that there exists a zero
vector. For this one, let's go through each axiom. Let $u =
\langle x_1, -x_1 \rangle$ and $v = \langle x_2, -x_2 \rangle$.
Notice that $u \oplus v = \langle x_1 + x_2, -x_1 - x_2 \rangle
\in \mathbb{R}^2$ and satisfies the condition $x_1 + x_2 + (-x_1
-x_2) = 0$. Therefore, A1 is satisfied. Moreover, consider the
following equations for A2 and A3 where $w = \langle x_3, -x_3
\rangle$.
\begin{align*}
    u \oplus v
    &= \langle x_1, -x_1 \rangle + \langle x_2, -x_2 \rangle
    = \langle x_1 + x_2, -x_1 - x_2 \rangle \\
    &= \langle x_2 + x_1, -x_2 - x_1 \rangle
    = \langle x_2, -x_2 \rangle + \langle x_1, -x_1 \rangle \\
    &= v \oplus u \\
    u \oplus (v \oplus w)
    &= \langle x_1, -x_1 \rangle + \left(
        \langle x_2, -x_2 \rangle + \langle x_3, -x_3 \rangle
    \right) \\
    &= \langle x_1, -x_1 \rangle +
        \langle x_2 + x_3, -x_2 - x_3 \rangle \\
    &= \langle x_1 + x_2 + x_3, -x_1 - x_2 - x_3 \rangle \\
    &= \langle x_1 + x_2, -x_1 - x_2 \rangle +
        \langle x_3, -x_3 \rangle \\
    &= (u \oplus v) \oplus w
\end{align*}
Therefore, A2 and A3 are satisfied. Furthermore, notice that
$\langle 0, 0 \rangle \in \mathbb{R}^2$ and $0 + 0 = 0$, satisfying
A4. Lastly for vector addition A5 holds as $\langle -x, x \rangle$
is the additive inverse of $\langle x, -x \rangle$ that satisfies
the provided conditions. Thus, all axioms for vector addition are met.

Continuing with scalar multiplication, notice that S1 and S3 satisfy
as $\langle x, y \rangle \in \mathbb{R}^2$ and $1 \odot \langle x, y
\rangle = \langle x, y \rangle$. Moreover, consider the following
equations for $\alpha, \beta \in \mathbb{R}$.
\begin{align*}
    \alpha \odot (\beta \odot u)
    &= \alpha \odot \langle \beta x_1, -\beta x_1 \rangle
    = \langle \alpha\beta x_1, -\alpha\beta x_1 \rangle
    = \alpha\beta \odot \langle x_1, -x_1 \rangle \\
    &= (\alpha\beta) \odot u \\
    (\alpha + \beta) \odot u
    &= \langle (\alpha + \beta) x_1, -(\alpha + \beta) x_1 \rangle
    = \langle \alpha x_1 + \beta x_1,
        -\alpha x_1 - \beta x_1 \rangle \\
    &= \langle \alpha x_1, -\alpha x_1 \rangle +
        \langle \beta x_1, -\beta x_1 \rangle
    = \alpha \langle x_1, -x_1 \rangle +
        \beta \langle x_1, -x_1 \rangle \\
    &= \alpha \odot u \oplus \beta \odot u \\
    \alpha \odot (u \oplus v)
    &= \alpha \langle x_1 + x_2, -x_1 - x_2 \rangle
    = \langle \alpha x_1 + \alpha x_2,
        -\alpha x_1 - \alpha x_2 \rangle \\
    &= \langle \alpha x_1, -\alpha x_1 \rangle +
        \langle \alpha x_2, -\alpha x_2 \rangle
    = \alpha \langle x_1, -x_1 \rangle +
        \alpha \langle x_2, -x_2 \rangle \\
    &= \alpha \odot u \oplus \alpha \odot v
\end{align*}
Thus, S2, S4, and S5 are all satisfied. Because all ten axioms for
a vector space are met, the set $\mathbb{S} = \{ \langle x, y \rangle
\in \mathbb{R}^2 \mid x + y = 0 \}$ is a vector space.
\end{solution}

\begin{problem}{}{}
Consider vectors in the form $\langle x, y, z \rangle$ where the
vector addition is defined as $\langle x, y, z \rangle \oplus
\langle x', y', z' \rangle = \langle x + x', y + y', z + z' \rangle$
and scalar multiplication as $\alpha \langle x, y, z \rangle =
\langle x^\alpha, y^\alpha, z^\alpha \rangle$. Identify if the vectors
form a vector space.
\end{problem}
\begin{solution}
First, notice that the vector addition adheres to the standard vector
addition. Therefore, we only have to look at scalar multiplication
if any axioms are violated. From a glance, S1, S2, and S3 appears to
be satisfied. Therefore, let's take a look at S4.
\[
(\alpha + \beta) \odot u
= \left(
    x^{\alpha + \beta}, y^{\alpha + \beta}, z^{\alpha + \beta}
\right)
\]
Notice that in general, $x^{\alpha + \beta} \neq x^\alpha + x^\beta$.
Therefore, the set of vectors does not form a vector space.
\end{solution}

Just as we have vector spaces $\mathbb{R}^n$, we could also define
spaces depending on the objects.

\begin{definition}{}{}
The symbol $\mathbb{P}^n$ denotes the vector space composed of
polynomials of degree at most $n$ and adheres to the ten properties.
$\mathbb{M}_{m \times n}$ represents the vector space with matrices
of order $m \times n$ that adheres to the ten properties.
\end{definition}

Therefore, we can say that $\mathbb{R}^n$ is more or less the same as
$\mathbb{M}_{1 \times n}$. With the understandings above, we can
derive some important theorems. Although most of them are intuitive,
it's always nice to formally prove them so let's do that now.

\begin{theorem}{}{Uniqueness of the Zero Vector}
For all vector spaces, zero vector, or the additive identity element,
is unique.
\end{theorem}
\begin{proof}
For the sake of contradiction, let $0_1, 0_2 \in V$ be distinct
additive identity element for a vector space $V$. By definition,
the following equations hold.
\begin{align*}
    0_1 \oplus 0_2 &= 0_1 \\
    0_1 \oplus 0_2 &= 0_2 \\
    \therefore 0_1 \oplus 0_2 &= 0_1 = 0_2
\end{align*}
Because $0_1 = 0_2$ is contradictory from the fact that $0_1$ and
$0_2$ are distinct, the initial assumption that there exists distinct
additive identity element in a vector space is false.
\end{proof}

Continuing from the uniqueness of the zero vector in a vector space,
we can assert the following theorem.

\begin{theorem}{}{Uniqueness of Additive Inverse}
For all $u \in V$ for any vector space $V$, $v \in V$ such that
$u \oplus v = 0$ is unique.
\end{theorem}
\begin{proof}
For the sake of contradiction, assume $v, w \in V$ are distinct
additive inverses of $u$. By definition and properties of vector
spaces, the following equations hold.
\[
v
= v \oplus 0
= v \oplus (u \oplus w)
= (v \oplus u) \oplus w
= 0 \oplus w
= w
\]
Because $v = w$ contradicts the definition that $v$ and $w$ are
unique, the initial assumption that there exists more than one
additive inverse of a vector in vector spaces is false.
\end{proof}

Just as $x = 0$ if $x + x = x$, we could do similar things in vector
spaces.

\begin{theorem}{}{Property of Zero Vectors I}
For all $v \in V$ and any vector space $V$, $v \oplus v = v$ if and
only if $v = 0$.
\end{theorem}
\begin{proof}
First, the first half of the statement can be proven. Consider the
following equations.
\begin{align*}
    (v \oplus v) \oplus (-v)
    &= v \oplus (-v) = 0 \\
    &= v \oplus (v \oplus (-v)) = v \oplus 0 = v
\end{align*}
Therefore, if $v \oplus v = v$, then $v = 0$. Moreover, if $v = 0$,
then $v \oplus v = 0 \oplus 0 = 0$. Therefore, the premise holds.
\end{proof}

Now let's take a look at theorems for scalar multiplications.

\begin{theorem}{}{Properties of Scalar Multiplications}
The following properties hold for all vector space $V$ over
$\mathbb{F}$, $u, \mathbf{0} \in V$, and any scalars $\alpha$.
\begin{enumerate}[noitemsep]
\item $0 \odot u = \mathbf{0}$
\item $\alpha \odot \mathbf{0} = \mathbf{0}$.
\item $(-1) \odot u = -u$
\item $u = -(-u)$
\item $\alpha \odot u = \mathbf{0}$ if and only if $\alpha = 0$ or
$u = \mathbf{0}$.
\end{enumerate}
\end{theorem}

Let's start with the first proof.

\begin{proof}
Consider the following equation.
\[
0 \odot u
= (0 + 0) \odot u
= 0 \odot u \oplus 0 \odot u
\]
By Theorem \ref{theo:Property of Zero Vectors I}, $0 \odot u =
\mathbf{0}$.
\end{proof}

Below is the proof for the second property.

\begin{proof}
Consider the following equation.
\[
\alpha \odot \mathbf{0}
= \alpha \odot (\mathbf{0} \oplus \mathbf{0})
= \alpha \odot \mathbf{0} \oplus \alpha \odot \mathbf{0}
\]
By Theorem \ref{theo:Property of Zero Vectors I}, $\alpha \odot
\mathbf{0} = \mathbf{0}$
\end{proof}

Continuing, here's the proof for the third one.

\begin{proof}
By the properties of vector spaces, $\mathbf{0} = 0 \odot u =
(1 - 1) \odot u = 1 \odot u \oplus (-1) \odot u = u \oplus (-1)
\odot u $. Because $(-1) \odot u$ is the additive inverse of $u$,
$(-1) \odot u = -u$.
\end{proof}

Below is the proof for the fourth one.

\begin{proof}
Notice that $-(-u)$ is the additive inverse of $-u$. However, by
definition, $u$ is the additive inverse of $-u$. By Theorem
\ref{theo:Uniqueness of Additive Inverse}, $u = -(-u)$.
\end{proof}

Finally, let's prove the last property.

\begin{proof}
It suffices to show the latter half of the statement as the first
half is shown by the first two properties. Consider the following
equations for $\alpha \neq 0$.
\begin{align*}
    \frac{1}{\alpha} \odot (\alpha \odot u)
    &= \frac{1}{\alpha} \odot \mathbf{0}
    = \mathbf{0} \\
    &= \left( \frac{1}{\alpha} \cdot \alpha \right) \odot u
    = 1 \odot u
    = u
\end{align*}
Therefore, if $\alpha \neq 0$, then $u = \mathbf{0}$. Similarly,
if $u \neq \mathbf{0}$, then $\alpha = 0$ as the product of nonzero
numbers is a nonzero number.
\end{proof}

We have covered the fundamental concepts and properties in vector
spaces. Before we move on with the next topic in vector spaces,
I wanted to discuss briefly on dot product as it will frequently
appear when you deal with vectors.

\begin{definition}{}{}
\textbf{Dot product} is an algebraic operation defined as the
following for vectors $a = [ a_1, \ldots, a_n]$ and $b =
[ b_1, \ldots, b_n]$
\[
a \cdot b
= \sum_{k = 1}^n a_k b_k
\]
\end{definition}

Here notice that the number of elements of the vectors in dot product
must be equal for the product to be defined. Below are some examples
of dot products.
\begin{align*}
\begin{bmatrix}
    1 \\
    2 \\
    3
\end{bmatrix}
\cdot
\begin{bmatrix}
    -1 \\
    -2 \\
    -3
\end{bmatrix}
&= 1(-1) + 2(-2) + 3(-3)
= -14 \\
\begin{bmatrix}
    0 & 1 & 2
\end{bmatrix}
\cdot
\begin{bmatrix}
    0 & 1 & 0
\end{bmatrix}
&= 0 \cdot 0 + 1 \cdot 1 + 2 \cdot 0
= 1
\end{align*}

Note that this is very different from the matrix multiplication
that we will discuss in the next note and we must write the
notation $\cdot$ to ensure that we are not doing matrix
multiplication. With that in mind, let's continue our discussion
with subspaces.

\end{document}


